%% 第26回編集委員会（物理第4巻）：組 電W8b
%% p4-00.tex  電-0　序章　発電所の電気は、どうやって家まで届く？（★1〜2。3段にしない）
%% 中身は計画書 tools/phys4-plan.md の B「電-0」。本文の章（電-1〜15）を読んで書いた、この巻の地図。第3巻の p3-00.tex・第2巻の p2-00.tex と同じ形。
%% 序章なので、じぶんでとこう・章末問題はない（例題は7段の見本3つ）。「この章で導く中学・高校の公式」の欄は置かない
%%   （P=VI・ジュール熱の持ち主は電-4 ⑤、送電の損失 R(P/V)^2 は電-11 ⑥、変圧器は電-12 ⑤。第19回 6-1-2・K-8）。
%%   はてな？の答え（損失 R(P/V)^2）は電-11 ⑥が持ち主なので、この章では「一般の式から、よく使う式へ」（phflow）を置かず、
%%   P=VI と RI^2（電-4 ⑤）に数を入れた見こみ（例題3・図）だけを見せて、式の形は電-11 ⑥・例題5 を指した。
%% 記号カードは置かない（電流 I・抵抗 R・電力 P・起電力 ℰ のカードは電-4 にある。U-6・X-2）。
%% 記号（この章）：電圧 V〔V〕、電流 I〔A〕、抵抗 R〔Ω〕、電力 P〔W〕、損失 P_損〔W〕、時間 t、距離 L、光の速さ c（この章では c を光の速さの意味だけで使う）。
%%   この章では回路の章（電-9〜13）の e・E・v・i は使わず、高校と同じ V・I で書く（Y-3。taisetsu で一言）。
%% 例題の並び：計画書の例題1（100 MW）・2（ドライヤー）・3（400 km）を、7段の見本の順に、例題1 ドライヤー・例題2 400 km・例題3 100 MW（はてな？）とした。
%% 数値はすべて Python で計算した：checks/物理第4巻/電-0/calc_p4_00.py（出力 calc_p4_00_out.txt）。
%% 本文の章とそろえた数値（同じ題材は、持ち主の章の値）：
%%   ・ドライヤー 100 V・1200 W・12 A・約8.3 Ω：電-4 例題1・⑤、電-10 例題2、電-14 例題5。電圧の変動 95〜107 V も電-4 例題1 と同じ。
%%   ・15 A の延長コード：電-4 ⑫・例題12（「〜とします」）。
%%   ・電子の速さ 約0.44 mm/s（12 A・2.0 mm²の銅線）：電-4 例題2。「知らせ」は光の速さに近い速さで伝わる：電-4 ②。
%%   ・架空送電線の上の速さはほぼ真空の光の速さ、50 Hz の波長 6000 km：電-13 ⑩。ケーブルの中は約3分の2：電-13 ①。
%%   ・100 MW・往復の抵抗 10 Ω・100 kV で 1000 A・10 MW（10 %）、500 kV で 200 A・0.4 MW（0.4 %）、275 kV で約1.3 %：電-11 例題5・じぶんでとこう。
%%   ・送電・配電の電圧（500・275・154・66 kV、6.6 kV、柱上変圧器 6600 V→210 V・105 V、単相3線式）・500 kV の送電は1973年から：電-11 ①。
%%   ・同期発電機 50 Hz で毎分3000回転（2極）：電-12 ⑧。変圧器の巻数の比：電-12 ⑤。発電機の起電力：電-8 ⑤。
%%   ・50 Hz・60 Hz、100 V の山 141 V：電-10 ①②。LED 電球 約0.1 A：電-4 の桁の感覚。10 W の LED 電球：数学 5-12⑥。
%%   ・火力発電所の効率 約4〜6割：熱-0（石炭 約0.40、コンバインドサイクル 約0.60）。
%%   ・電-16 はまだ書かれていないので、計画書 B の節の番号で指した（② 周波数と需給、⑦ 電気をためる）。
%%   ・各章のはてな？の答えの節は、書き上がった原稿で確かめた（電-1 ②、電-2 ④、電-3 ③、電-4 ⑤、電-5 ①、電-6 ⑤、電-7 ②、電-8 ⑫、
%%     電-9 ⑤、電-10 ⑦、電-11 ⑧、電-12 ⑧、電-13 ⑧、電-14 ⑨、電-15 ⑩）。電-16 は計画書の②。
%% 出典（第26回 手引き 2章の16：書くときに確かめた。11月5日）
%%   ・家庭の電圧：電気事業法施行規則 第38条第1項（標準電圧 100 V は「百一ボルトの上下六ボルトを超えない値」、200 V は「二百二ボルトの上下二十ボルトを
%%       超えない値」）、同第2項（周波数は標準周波数に等しい値）。e-Gov 法令 API https://laws.e-gov.go.jp/api/1/lawdata/407M50000400077 で本文を読んだ。
%%   ・送電・配電の電圧の段（Y-7：二次資料を2つ以上で照らし合わせた。本文は「2020年代の資料。会社や地域で違う」と書いた）：
%%       関西電力「電気事業のあらまし」2011・2012 の送電の頁（発電所で 27万5,000〜50万 V に上げ、一次変電所で 15.4万〜7.7万 V に下げる）、
%%       中小機構 J-Net21 の省エネ Q&A（Q1228。配電用変電所から柱上変圧器までは 6,600 V、柱上変圧器から家庭までは 100/200 V）、
%%       パナソニックの用語解説「柱上変圧器」（6,600 V を 100〜200 V に下げる）、コトバンク「送電」（日本大百科全書。公称電圧 500・275・220・187・154・
%%       110・77・66 kV など）、日本製鉄の季刊誌（2020年）。いずれも WebSearch の抜粋で読んだ（kepco.co.jp・panasonic.biz・cgworld.jp・tepco.co.jp・
%%       energy-kyoiku.meti.go.jp は作業環境から届かなかった）。電-11 の頭のコメントの出典（2025年の時点）とも同じ値。
%%   ・電気をたくさん・長くためる方法の研究（まだわかっていないこと）：アメリカのエネルギー省（DOE）の Long Duration Storage Shot（2021年7月。
%%       10時間をこえて電気をためる大きな蓄電の費用を、10年で9割下げる目標。2025年1月のファクトシート）、Sandia 国立研究所の記事（2023年）、
%%       Utility Dive の記事（2021年）。WebSearch の抜粋で照らし合わせた（energy.gov は作業環境から届かなかった）。
%%       揚水発電がいちばん大きな長時間の蓄電であることも DOE の説明による。本文は数を出さず「研究が続いている」とだけ書いた。
%%   ・〔要確認〕なし（Wikipedia は出典にしていない）。
\makeatletter
% 会話の1行  \sxatalk{y}{sora|mina|delta}{名前}{せりふ}（p3-00 の \rxatalk と同じ形）
\newcommand\sxatalk[4]{%
  \bkChar{#2}{(7,#1-6.5)}{12mm}
  \node[rectangle callout, callout absolute pointer={(15.5,#1)}, anchor=west, draw=bkTheme!80!black, line width=0.8pt,
    fill=white, rounded corners=2.5mm, inner xsep=3mm, inner ysep=1.8mm, text width=103mm, align=left,
    font=\small] at (21,#1) {#4};
  \node[anchor=north, font=\sffamily\footnotesize, text=bkGray, inner sep=0pt] at (7,#1-7.5) {#3};}
% 送電線の鉄塔  \sxatower{(足のまん中)}{高さ}
\newcommand\sxatower[2]{\begin{scope}[shift={#1}]
  \draw[bkGray!80!black, line width=0.7pt] (-5,0) -- (-1.2,#2) -- (1.2,#2) -- (5,0);
  \foreach \k in {1,...,5} {\draw[bkGray!80!black, line width=0.35pt] ({-5+3.8*(\k-1)/5},{#2*(\k-1)/5}) -- ({5-3.8*\k/5},{#2*\k/5})
     ({5-3.8*(\k-1)/5},{#2*(\k-1)/5}) -- ({-5+3.8*\k/5},{#2*\k/5});}
  \draw[bkGray!80!black, line width=0.8pt] (-6,{#2*0.72}) -- (6,{#2*0.72}) (-6,{#2*0.88}) -- (6,{#2*0.88});
\end{scope}}
% 家  \sxahouse{(左下)}{色}
\newcommand\sxahouse[2]{\begin{scope}[shift={#1}]
  \filldraw[fill=#2!12, draw=#2!70!black, line width=0.7pt] (0,0) rectangle (12,9);
  \filldraw[fill=#2!35, draw=#2!70!black, line width=0.7pt] (-1.5,9) -- (6,15) -- (13.5,9) -- cycle;
  \filldraw[fill=white, draw=#2!70!black, line width=0.5pt] (4.5,0) rectangle (7.5,5);
\end{scope}}
% 電柱と柱上変圧器  \sxapole{(根もと)}{高さ}
\newcommand\sxapole[2]{\begin{scope}[shift={#1}]
  \fill[bkGray!60] (-0.7,0) rectangle (0.7,#2);
  \draw[bkGray!80!black, line width=0.8pt] (-4,{#2-1}) -- (4,{#2-1});
  \filldraw[fill=bkGray!25, draw=bkInk, line width=0.5pt] (0.8,{#2-8}) rectangle (4.2,{#2-3.5});
\end{scope}}
\makeatother
\tikzset{sxa lab/.style={font=\sffamily\footnotesize, inner sep=1pt, text=bkInk},
  sxa arr/.style={-{Stealth[length=2.6mm, width=2.2mm]}, line width=1.1pt},
  sxa box/.style={bk node, minimum width=28mm, text width=26.5mm, minimum height=20mm, font=\sffamily\footnotesize, align=center, inner sep=1pt},
  sxa dim/.style={bkGray!80!black, line width=0.5pt, {Stealth[length=1.8mm]}-{Stealth[length=1.8mm]}}}

\phchapter{0}{発電所の電気は、\\どうやって家まで届く？}{送電線は、なぜとても高い電圧で電気を送る？}{\bklevelrange{1}{2}}{発電所の電気は、どうやって家まで届く？}
\bkmichishirube{2-10,2-15,5-12,力-4,熱-0}{電-4,電-8,電-11,電-12,電-13}

\begin{goal}
\begin{itemize}[label={\color{bkTheme}●}, leftmargin=1.6\zw, itemsep=0pt, topsep=1pt]
\item \bklevel{1}\enspace 発電所から家のドライヤーまで、電気が通る道すじと、とちゅうで働く機械（発電機・変圧器・送電線）を説明できる。
\item \bklevel{2}\enspace 何を無視したかを書いてから、電流・電力・電線で失われる電力を、7つの段で求められる。
\item \bklevel{1}\enspace 身近な「なぜ？」の答えが、この巻のどの章にあるかを探せる。
\end{itemize}
\end{goal}

\begin{hatena}
\noindent\begin{minipage}[c]{0.54\linewidth}\raggedright
山の中や海辺には、大きな鉄塔が並び、太い電線が張られています。この電線（送電線）には、何十万ボルトという、とても高い電圧がかかっています。一方、家のコンセントの電圧は\qty{100}{V}です。\par\vspace{3pt}
高い電圧は危ないのに、なぜわざわざ高くして送るのでしょう。家で使う\qty{100}{V}のまま送ったら、何が困るのでしょう。
\end{minipage}\hfill
\begin{minipage}[c]{0.43\linewidth}\centering
\begin{tikzpicture}[x=0.82mm, y=0.86mm]
  \fill[bkGray!25] (-2,-2) rectangle (60,0);
  \draw[bkInk, line width=0.6pt] (-2,0) -- (60,0);
  \sxatower{(8,0)}{36}
  \draw[bkInk, line width=0.6pt] (2,25.9) .. controls (20,17) and (32,17) .. (44.4,21)
    (14,25.9) .. controls (28,20) and (36,20) .. (44.4,23) (2,31.7) .. controls (22,24) and (34,25) .. (44.4,25);
  \sxapole{(44,0)}{26}
  \draw[bkInk, line width=0.5pt] (46.5,18) .. controls (50,14) and (52,13) .. (54,12);
  \sxahouse{(48,0)}{bkOrange}
  \node[sxa lab, anchor=south, fill=white] at (22,33) {何十万 V};
  \node[sxa lab, anchor=south] at (54,15.5) {\qty{100}{V}};
  \node[font=\sffamily\bfseries\Large, text=phHatena!80!black] at (30,8) {?};
\end{tikzpicture}
\end{minipage}
\end{hatena}

この本は、身の周りの「電気」と「磁石」の「なぜ？」に、物理で答える本です。序章では、発電所でつくった電気が、家のドライヤーに届くまでの旅を追いながら、この本で何を学ぶか、どの章で学ぶかを見ていきます。詳しい計算は、それぞれの章でします。いまはわからない式やことばが出てきても、だいじょうぶです。

\begin{yattemiyou}[電気の通り道を調べよう]
\textbf{用意するもの}：単3乾電池2本と電池ボックス、豆電球（ソケットつき）、導線、アルミはく・10円玉・クリップ・消しゴム・紙\par\vspace{2pt}
\noindent\begin{minipage}[c]{0.62\linewidth}
\begin{enumerate}[label={\bkcircled[bkGreen]{\arabic*}}, leftmargin=2.4\zw, itemsep=1pt, topsep=0pt]
\item 乾電池1本・豆電球・導線で、ぐるっと1つの輪をつくり、明かりをつける。導線を1か所外すと、どうなるだろうか。
\item 輪のとちゅうを1か所切り、その間に、アルミはく・10円玉・クリップ・消しゴム・紙を1つずつはさむ。明かりがつくのはどれだろうか。
\item 乾電池を2本、＋と−をつないで（直列に）入れると、明るさはどう変わるだろうか。
\end{enumerate}
\end{minipage}\hfill
\begin{minipage}[c]{0.35\linewidth}\centering
\begin{tikzpicture}[x=1mm, y=1mm]
  \draw[bkInk, line width=0.9pt] (0,9) -- (0,24) -- (32,24) -- (32,15.5) (32,8.5) -- (32,0) -- (0,0) -- (0,5.5);
  \draw[bkInk, line width=1.1pt] (-4.5,8.5) -- (4.5,8.5);
  \draw[bkInk, line width=2.2pt] (-2.2,5.5) -- (2.2,5.5);
  \node[sxa lab, anchor=east] at (-4.8,10.3) {$+$};
  \filldraw[fill=bkYellow!40, draw=bkInk, line width=0.8pt] (32,12) circle (3.4);
  \draw[bkInk, line width=0.6pt] (29.6,9.6) -- (34.4,14.4) (29.6,14.4) -- (34.4,9.6);
  \fill[white] (13,22) rectangle (19,26);
  \filldraw[fill=bkGray!20, draw=bkInk, line width=0.6pt] (13.5,22.5) rectangle (18.5,25.5);
  \node[sxa lab, anchor=south] at (16,26.3) {ここにはさむ};
  \node[sxa lab, anchor=west] at (5.5,7) {乾電池};
  \node[sxa lab, anchor=east] at (28,12) {豆電球};
\end{tikzpicture}
\end{minipage}\par\vspace{7pt}
\noindent\textbf{約束}：乾電池の＋と−を、豆電球を入れずに導線だけでつながない（導線と電池が熱くなる）。家のコンセントでは、決して試さない。
\end{yattemiyou}

\noindent やってみようの答え：①導線を1か所でも外すと、明かりは消えます。電気は、乾電池の＋から出て、豆電球を通り、−へもどる\textbf{1つの輪}がつながったときだけ流れます。この輪を\textbf{回路}といいます。\bkidx{かいろ}{回路}②アルミはく・10円玉・クリップのような金属をはさむと明かりがつき、消しゴムや紙ではつきません。金属は電気をよく通し、ゴムや紙はほとんど通しません（電-4 ①）。③乾電池2本を直列にすると、電気を押し出すはたらき（電圧）が2倍になり、豆電球は明るくなります（電-5 ②）。

発電所から家のドライヤーまでも、じつは同じです。行きの線と帰りの線で、何百 km もの大きな輪がつながっていて、発電所の発電機が乾電池の役をしています。乾電池と違うのは、電気を押す向きが1秒に何十回も入れかわること（交流）と、とちゅうで電圧を何度も変えることです（図\ref{ph:電-0:fig-wa}）。
\begin{bkfigure}
\begin{tikzpicture}[x=1mm, y=1mm]
  % (a) 乾電池と豆電球
  \draw[bkInk, line width=0.9pt] (0,9) -- (0,16) -- (36,16) -- (36,11.4) (36,4.6) -- (36,0) -- (0,0) -- (0,5.5);
  \draw[bkInk, line width=1.1pt] (-4.5,8.5) -- (4.5,8.5);
  \draw[bkInk, line width=2.2pt] (-2.2,5.5) -- (2.2,5.5);
  \node[sxa lab, anchor=east] at (-4.8,10.3) {$+$};
  \filldraw[fill=bkYellow!40, draw=bkInk, line width=0.8pt] (36,8) circle (3.4);
  \draw[bkInk, line width=0.6pt] (33.6,5.6) -- (38.4,10.4) (33.6,10.4) -- (38.4,5.6);
  \draw[sxa arr, bkBlue!80!black] (13,16) -- (23,16);
  \draw[sxa arr, bkBlue!80!black] (23,0) -- (13,0);
  \node[sxa lab, anchor=west] at (5.5,7) {乾電池};
  \node[sxa lab, anchor=east] at (32,8) {豆電球};
  \node[sxa lab, anchor=north] at (18,-2.5) {(a)\enspace 机の上の回路（数十 cm）};
  % (b) 発電機とドライヤー
  \filldraw[fill=bkOrange!15, draw=bkInk, line width=0.8pt] (62,8) circle (5.5);
  \node[sxa lab] at (62,8) {発電機};
  \draw[bkInk, line width=0.9pt] (62,13.5) -- (62,16) -- (118,16) -- (118,13) (118,3) -- (118,0) -- (62,0) -- (62,2.5);
  \filldraw[fill=bkRed!15, draw=bkInk, line width=0.8pt] (108,3) rectangle (128,13);
  \node[sxa lab] at (118,8) {ドライヤー};
  \draw[sxa arr, bkBlue!80!black] (72,16) -- (82,16);
  \draw[sxa arr, bkBlue!80!black] (82,0) -- (72,0);
  \node[sxa lab, anchor=south] at (97,16.6) {行きの線};
  \node[sxa lab, anchor=south] at (93,0.6) {帰りの線};
  \node[sxa lab, anchor=north] at (93,-2.5) {(b)\enspace 発電所から家まで（何百 km）};
\end{tikzpicture}
\bkfigcaption{どちらも1つの輪（回路）。(b) では変圧器を省いた。交流では、電流の向き（矢印）が1秒に何十回も入れかわる。}\label{ph:電-0:fig-wa}
\end{bkfigure}

\begin{omoidasou}
\bkrecall{2-15}{電圧（電流を押し出すはたらき）の単位は V（ボルト）、電流の単位は A（アンペア）。抵抗を流れる電流は、電圧に比例する（オームの法則。2-15④）。}
\bkrecall{2-10}{単位の前の k（キロ）は1000倍、M（メガ）は100万倍（2-10③）。\qty{500}{kV}は\qty{500000}{V}、\qty{100}{MW}は1億 W。}
\bkrecall{5-12}{電力〔W〕は、1秒あたりに使うエネルギー。電力×時間が電力量で、$\qty{1}{kWh}=\qty{3.6e6}{J}$（5-12⑥）。}
\bkrecall{力-4}{仕事率は1秒あたりの仕事で、単位は W（力-4 ⑧）。電気の電力も、同じ量。}
\bkrecall{熱-0}{火力発電所では、燃料の熱のうち、約4割（石炭）から約6割（コンバインドサイクル）が電気になる。}
\end{omoidasou}

%% ================================================================
\Needspace{16\baselineskip}
\bksection{hon}{この巻の地図}

\bksubsection{1}{発電所から家まで：電気の旅}
\leavevmode\label{ph:電-0:sec1}発電所でつくった電気は、図\ref{ph:電-0:fig-tabi}のように、何回も電圧を変えながら、家のコンセントまで届きます。旅の順に、どこで、どの章の物理が働いているかを見ていきましょう。
\begin{bkfigure}
\begin{tikzpicture}[x=1mm, y=1mm]
  \node[sxa box] (a1) at (14,0) {\bkcircled[bkBlue]{1}\,\textbf{発電機}\\磁石とコイルで\\交流をつくる\\電-8 ⑤};
  \node[sxa box] (a2) at (48,0) {\bkcircled[bkBlue]{2}\,\textbf{変圧器}\\電圧を上げる\\（何十万 V へ）\\電-12 ⑤};
  \node[sxa box] (a3) at (82,0) {\bkcircled[bkBlue]{3}\,\textbf{送電線}\\500・275 kV\\三相の3本1組\\電-11・電-13};
  \node[sxa box] (a4) at (116,0) {\bkcircled[bkBlue]{4}\,\textbf{変電所}\\154・66 kV に\\段々に下げる\\電-12 ⑤};
  \node[sxa box] (a5) at (116,-29) {\bkcircled[bkBlue]{5}\,\textbf{配電用変電所}\\6.6 kV に下げ\\電柱の線へ\\電-11 ①};
  \node[sxa box] (a6) at (82,-29) {\bkcircled[bkBlue]{6}\,\textbf{柱上変圧器}\\6600 V を\\100・200 V に\\電-11 ①};
  \node[sxa box] (a7) at (48,-29) {\bkcircled[bkBlue]{7}\,\textbf{コンセント}\\\qty{100}{V}の交流\\50・60 Hz\\電-10 ①};
  \node[sxa box] (a8) at (14,-29) {\bkcircled[bkBlue]{8}\,\textbf{ドライヤー}\\\qty{1200}{W}\\電気が熱になる\\電-4 ⑤};
  \draw[bk link] (a1) -- (a2); \draw[bk link] (a2) -- (a3); \draw[bk link] (a3) -- (a4);
  \draw[bk link] (a4) -- (a5); \draw[bk link] (a5) -- (a6); \draw[bk link] (a6) -- (a7); \draw[bk link] (a7) -- (a8);
\end{tikzpicture}
\bkfigcaption{発電所から家のドライヤーまで、電気の旅。電圧の段は一例で、会社や地域で違う。どの段にも、その物理を学ぶ章を書いた。}\label{ph:電-0:fig-tabi}
\end{bkfigure}
\begin{enumerate}[label={\bkcircled[bkBlue]{\arabic*}}, leftmargin=2.4\zw, itemsep=2pt, topsep=2pt]
\item \textbf{発電機}　火力発電所では、燃料の熱で水蒸気をつくり、タービンという羽根車を回します（熱-0）。タービンにつながった\textbf{発電機}は、磁石とコイルを回して電気をつくる機械です。\bkidx{はつでんき}{発電機}コイルの中を通る磁場が変わると電圧が生まれる\textbf{電磁誘導}（電-8 ③）を使うので、生まれる電圧は自然に、＋と−を行ったり来たりする\textbf{交流}になります（電-8 ⑤）。\bkidx{こうりゅう}{交流}東日本（\qty{50}{Hz}）の火力発電所の発電機の多くは、1分間に3000回、1秒に50回まわります。発電機の回る速さが、そのまま交流の周波数になるのです（電-12 ⑧）。
\item \textbf{変圧器で上げる}　大きな発電所では、発電所を出る所で、\textbf{変圧器}が電圧を何十万 V に上げます。\bkidx{へんあつき}{変圧器}変圧器は、鉄の輪に2つのコイルを巻いただけの機械で、2つのコイルの巻数の比で電圧を変えます（電-12 ⑤）。動く部分がなく、交流でだけ働きます。
\item \textbf{送電線}　鉄塔の上の電線（\textbf{送電線}）で、遠くの町まで送ります。\bkidx{そうでんせん}{送電線}日本の送電線の電圧には、\qty{500}{kV}・\qty{275}{kV}・\qty{154}{kV}・\qty{66}{kV}などの段があります（2020年代の電力会社などの資料。会社や地域で違います）。\qty{500}{kV}で送るようになったのは1973年です（電-11 ①）。鉄塔の電線が3本ずつ組になっているのは、\textbf{三相交流}という送り方をしているからです（電-11 のはてな？）。何百 km もの長い線になると、線そのものがコイルやコンデンサーのように働くことも考えます（電-13 ⑩）。
\item \textbf{変電所で下げる}　町に近づくと、\textbf{変電所}の変圧器で、\qty{154}{kV}、\qty{66}{kV}と、段々に電圧を下げます。\bkidx{へんでんしょ}{変電所}
\item \textbf{配電線}　\textbf{配電用変電所}で\qty{6.6}{kV}（\qty{6600}{V}）に下げ、町の電柱の上の3本の線（\textbf{配電線}）で配ります。\bkidx{はいでんせん}{配電線}電柱の上の電線にとまった鳥が感電しないわけは、電-5 のはてな？で調べます。
\item \textbf{柱上変圧器}　電柱の上の箱（\textbf{柱上変圧器}）が、\qty{6600}{V}を\qty{105}{V}と\qty{210}{V}に下げ、3本の線で家に引きこみます。\bkidx{ちゅうじょうへんあつき}{柱上変圧器}家では、この3本の線から\qty{100}{V}と\qty{200}{V}を取り出して使います（単相3線式。電-11 ①）。
\item \textbf{コンセント}　家に届く電圧は、法律（電気事業法施行規則）で、\qty{100}{V}は「\qty{101}{V}の上下\qty{6}{V}をこえない値」（\qtyrange{95}{107}{V}）、\qty{200}{V}は「\qty{202}{V}の上下\qty{20}{V}をこえない値」（\qtyrange{182}{222}{V}）に保つと決められています。コンセントの電気は交流で、東日本は\qty{50}{Hz}、西日本は\qty{60}{Hz}です（電-10 ①）。
\item \textbf{ドライヤー}　ドライヤーの中の細い金属の線（電熱線）に電流が流れると、線が熱くなり、風を温めます（電-4 ⑤）。発電所でつくった電気のエネルギーが、ここで熱になって、旅が終わります。
\end{enumerate}
\noindent この旅には、2つの見方があります。1つは、電線の中を流れる電流と、電線の間の電圧で考える\textbf{回路}の見方です（第1部・第3部）。もう1つは、電線の周りの空間に広がる\textbf{電場}と\textbf{磁場}で考える\textbf{場}の見方です（第1部・第2部・第4部）。じつは、電気のエネルギーは、電線の中ではなく、電線の周りの空間を、電場と磁場にのって運ばれていることが、電-14 ⑧でわかります。この本では、2つの見方を行ったり来たりしながら学びます。

\bksubsection{2}{電気をはかる4つの量：電圧・電流・抵抗・電力}
\leavevmode\label{ph:電-0:sec2}電気の旅を数で追うには、表\ref{ph:電-0:tab-4ryo}の4つの量を使います。電気の流れは目に見えないので、はじめは水の流れにたとえるとわかりやすいでしょう。
\begin{bkcenter}
\small
\begin{tblr}{colspec={Q[l,m,wd=10mm] Q[l,m,wd=17mm] X[l,m] Q[l,m,wd=40mm] Q[l,m,wd=15mm]}, width=\linewidth, rowsep=1.1pt, colsep=3pt,
  hlines={0.4pt,bkLine}, row{1}={bg=bkTheme!12, font=\sffamily\bfseries\footnotesize}}
量 & 記号〔単位〕 & 水の流れにたとえると & 身近な値 & 学ぶ章 \\
電圧 & {$V$〔V〕\\ボルト} & 水を押し出す、高さの差 & {乾電池\qty{1.5}{V}\\コンセント\qty{100}{V}} & 電-2・電-4 \\
電流 & {$I$〔A〕\\アンペア} & 1秒あたりに流れる水の量 & {LED 電球 約\qty{0.1}{A}\\ドライヤー\qty{12}{A}} & 電-4 \\
抵抗 & {$R$〔\si{\ohm}〕\\オーム} & 水路の流れにくさ（細さ・長さ） & {ドライヤーの電熱線\\約\qty{8.3}{\ohm}} & 電-4 \\
電力 & {$P$〔W〕\\ワット} & 水が1秒あたりにする仕事 & {LED 電球 約\qty{10}{W}\\ドライヤー\qty{1200}{W}} & 電-4 \\
\end{tblr}
\bktabcaption{電気をはかる4つの量。水のたとえは目安で、合わない所もある（電気は、線の周りの電場と磁場がエネルギーを運ぶ：電-14 ⑧）。}\label{ph:電-0:tab-4ryo}
\end{bkcenter}
4つの量の間には、中学理科で習う2つの関係があります。
\begin{itemize}[leftmargin=1.2\zw, itemsep=2pt, topsep=2pt, label={\color{bkTheme}●}]
\item \textbf{オームの法則}　$V=RI$。抵抗に流れる電流は、かけた電圧に比例する（2-15④）。なぜそうなるかは、電場が金属の中の電子を押すしくみから、電-4 ③で導きます。\bkidx{おーむのほうそく}{オームの法則}
\item \textbf{電力}　$P=VI$。1秒あたりに使われるエネルギーは、電圧と電流の積（電-4 ⑤で導く）。電圧が同じなら、電流が多いほど電力は大きい。\bkidx{でんりょく}{電力}
\end{itemize}
電圧の大きさは、乾電池から送電線まで、とても大きく違います（図\ref{ph:電-0:fig-volt}）。図の目もりは、1つ右へ進むごとに10倍になります（3-12 の対数の目もり）。送電線の\qty{500}{kV}は、乾電池の約30万倍、家のコンセントの5000倍です。
\begin{bkfigure}
\begin{tikzpicture}[x=1mm, y=1mm]
  % x = 5 + 20*log10(V)
  \draw[bkInk, line width=0.8pt, -{Stealth[length=2.4mm]}] (0,0) -- (131,0);
  \foreach \k in {0,...,6} {\draw[bkInk, line width=0.5pt] (5+\k*20,0) -- (5+\k*20,-1.6);}
  \foreach \x/\t in {5/1\,V, 25/10\,V, 45/100\,V, 65/1\,kV, 85/10\,kV, 105/100\,kV, 125/1\,MV} {\node[sxa lab, anchor=north] at (\x,-2.2) {\t};}
  \fill[bkGreen!60!black] (8.5,0) circle (1.3);
  \draw[bkOrange!90, line width=3pt] (45,2.2) -- (51.0,2.2);
  \fill[bkPurple] (81.4,0) circle (1.3);
  \draw[bkRed!70, line width=3pt] (101.4,2.2) -- (119.0,2.2);
  \node[sxa lab, align=center, anchor=south] at (8.5,3.5) {乾電池\\\qty{1.5}{V}};
  \node[sxa lab, align=center, anchor=south] at (48,4.2) {家\\\qty{100}{V}・\qty{200}{V}};
  \node[sxa lab, align=center, anchor=north] at (81.4,-7.5) {配電線\\\qty{6600}{V}};
  \node[sxa lab, align=center, anchor=south] at (110.2,4.2) {送電線\\\qty{66}{kV}〜\qty{500}{kV}};
\end{tikzpicture}
\bkfigcaption{電圧のものさし。目もりが1つ進むごとに10倍。丸は1つの値、太い線は範囲。}\label{ph:電-0:fig-volt}
\end{bkfigure}
\noindent コンセントの「\qty{100}{V}」は、交流の電圧を、同じはたらきをする直流の電圧で言いかえた値（\textbf{実効値}）です。電圧の山は約\qty{141}{V}あり、1秒に50回か60回、＋と−が入れかわります（電-10 ②）。この本で交流の電圧・電流を書くときは、ことわりがなければ実効値です。

%% ================================================================
\Needspace{12\baselineskip}
\bksubsection{3}{例題の解き方：7つの段}
\leavevmode\label{ph:電-0:sec3}この本の例題は、いつも同じ7つの段で解きます。①場面、②何を無視するか（理想化）、③図と記号、④使う式を選ぶ理由、⑤式を立てる、⑥解く、⑦答えを確かめる（単位・大きさ・極端な場合）です。数学の本の例題よりも段が多いのは、式を立てる前の「モデルを作る」部分と、答えを出したあとの「確かめる」部分を、ていねいに書くためです。何を無視したかは、式を立てる前に、次のような「モデル化の札」にも書きます。
\begin{phmodel}[ドライヤーを1本の電熱線と考える]
\phignore{風を送るファンのモーターが使う電力。コンセントの電圧の小さな変動。電熱線の温度による抵抗の変化。}
\phassume{\qty{100}{V}がそのまま電熱線にかかり、使われる電力が、全部、熱になる。}
\phlimit{スイッチを入れた直後で、電熱線がまだ冷たいとき（抵抗が小さく、電流が少し多い）。}
\end{phmodel}
\bkidx{もでる}{モデル}このように、知りたいことに関係の小さいものを思い切って無視して作った、簡単にした考え方を\textbf{モデル}といいます。見本を2つ見ておきましょう。

\begin{phreidai}[ドライヤーに流れる電流]\label{ph:電-0:ex1}
\qty{100}{V}・\qty{1200}{W}と書かれたドライヤーがあります。流れる電流を求めましょう。また、10分間使ったときの電力量は何 kWh ですか。
\phbamen
旅の終わりの⑧です。コンセントに届いた電気が、ドライヤーで熱になります。
\phmushi
モデル化の札のとおり、ファンの電力と、電圧の変動（\qtyrange{95}{107}{V}）を無視し、\qty{100}{V}・\qty{1200}{W}のまま使うとします。
\phzu
電圧 $V=\qty{100}{V}$、電力 $P=\qty{1200}{W}$、電流 $I$〔A〕、時間 $t=10$ 分$=\qty{600}{s}$。
\phriyu
電圧と電力がわかっているので、電力の式 $P=VI$（②、電-4 ⑤）から電流が出せます。電力量は「電力×時間」です（5-12⑥）。
\phtateru
\[I=\frac{P}{V},\qquad W=Pt\]
\phtoku
\[I=\frac{1200}{100}=\qty{12}{A}\]
\[W=1200\times600=\qty{720000}{J}=\frac{720000}{3.6\times10^6}\ \text{kWh}=\qty{0.20}{kWh}\]
\phtashikame
単位：W÷V＝(V·A)÷V＝A。\\
大きさ：LED 電球（約\qty{0.1}{A}）の120倍です。延長コードには、流してよい電流（たとえば\qty{15}{A}とします）が書いてあり、ドライヤー1台で、その8割を使います（電-4 例題12）。\\
極端な場合：電力が半分の\qty{600}{W}なら、電流も半分の\qty{6}{A}です。
\kotae \qty{12}{A}、\qty{0.20}{kWh}
\end{phreidai}
\noindent 電線の中の電子の動きを見ると、おどろくことがわかります。断面積\qty{2.0}{mm^2}の銅線を\qty{12}{A}が流れるとき、電子の流れの速さは、1秒に約\qty{0.44}{mm}しかありません（電-4 例題2）。それなのに、スイッチを入れると、ドライヤーはすぐ動きます。導線の中は、はじめから自由電子でぎっしりつまっているからです（図\ref{ph:電-0:fig-tama}）。
\begin{bkfigure}
\begin{tikzpicture}[x=1mm, y=1mm]
  \fill[bkGray!12] (10,0) rectangle (112,8);
  \draw[bkInk, line width=0.8pt] (10,0) -- (112,0) (10,8) -- (112,8);
  \foreach \k in {0,...,16} {\filldraw[fill=bkBlue!25, draw=bkBlue!70!black, line width=0.4pt] (13.2+6*\k,4) circle (2.8);}
  \filldraw[fill=bkBlue!25, draw=bkBlue!70!black, line width=0.4pt] (5,4) circle (2.8);
  \draw[sxa arr, bkRed!80!black] (-6,4) -- (1.5,4);
  \filldraw[fill=bkBlue!25, draw=bkBlue!70!black, line width=0.4pt] (117,4) circle (2.8);
  \draw[sxa arr, bkRed!80!black] (121,4) -- (128,4);
  \node[sxa lab, anchor=north] at (2,-1.5) {1つ押しこむ};
  \node[sxa lab, anchor=north] at (119,-1.5) {すぐ1つ出る};
  \node[sxa lab, anchor=south] at (61,9) {管の中は、はじめから玉でぎっしり（導線の中の自由電子）};
\end{tikzpicture}
\bkfigcaption{玉がつまった管。1つ押しこむと、反対のはしの玉がすぐ出る。「押せ」という知らせが速く伝わる。}\label{ph:電-0:fig-tama}
\end{bkfigure}

\begin{phreidai}[電気の「知らせ」は 400 km を何秒で進む？]\label{ph:電-0:ex2}
東京から大阪まで、長さ\qty{400}{km}の送電線があるとします。片方のはしで電圧を変えたとき、その「知らせ」がもう片方のはしに届くまでの時間を求めましょう。空気の中に張った送電線では、知らせは真空の光の速さ $c=\qty{3.00e8}{m/s}$ で進むとします。
\phbamen
電線の中は、はじめから自由電子でぎっしりつまっています。片はしで電圧を変えると、「押せ」という知らせが電線にそって伝わり、先の電子が動き始めます（電-4 ②）。
\phmushi
電線のとちゅうの変電所や、電線のくねくねした道すじを無視し、まっすぐな\qty{400}{km}の線とします。知らせの速さは、空気の中の電線では光の速さとほとんど同じなので（電-13 ⑩）、光の速さとします。
\phzu
長さ $L=\qty{400}{km}=\qty{4.00e5}{m}$、速さ $c=\qty{3.00e8}{m/s}$、時間 $t$〔s〕。
\phriyu
知らせは一定の速さで進むので、「時間＝距離÷速さ」が使えます。
\phtateru
\[t=\frac{L}{c}\]
\phtoku
\[t=\frac{4.00\times10^5}{3.00\times10^8}\approx\qty{1.33e-3}{s}\approx\qty{1.3}{ms}\]
\phtashikame
単位：m÷(m/s)＝s。\\
大きさ：1000分の1.3秒です。一方、電子そのものが1秒に約\qty{0.44}{mm}で\qty{400}{km}進もうとすると、$4.00\times10^5\div(4.4\times10^{-4})\approx9\times10^8$ 秒、約29年もかかります。届くのは電子ではなく、「押せ」という知らせなのです。しかも交流では、電子は1秒に50回か60回、向きを変えながら、その場で前後にゆれているだけです。\\
極端な場合：テレビのケーブルの中のように、線が樹脂でつつまれていると、知らせは光の速さの約3分の2になります（電-13 ①）。
\kotae 約\qty{1.3}{ms}
\end{phreidai}
\noindent \qty{50}{Hz}の交流が1回ゆれる時間（\qty{20}{ms}）にくらべると、\qty{1.3}{ms}はその約15分の1です。数百 km をこえる長い送電線では、このおくれが無視できなくなり、線の上を進む「電気の波」として考える必要が出てきます（電-13 ⑥⑩）。

%% ================================================================
\Needspace{14\baselineskip}
\bksubsection{4}{高い電圧で送るわけ}
\leavevmode\label{ph:電-0:sec4}いよいよ、\bkhatenaref を数で考えましょう。送電線も電線なので、小さいけれど抵抗があります。電流が流れると、ドライヤーの電熱線と同じように熱が出て、送った電力の一部が、とちゅうで失われてしまいます。この熱の電力は、電線の抵抗を $R$、電流を $I$ とすると $RI^2$（ジュール熱。電-4 ⑤）で、\textbf{電流の2乗}に比例します。\bkidx{そうでんのそんしつ}{送電の損失}

\begin{phreidai}[100 MW を高い電圧で送る（\bkhatenaref）]\label{ph:電-0:ex3}
発電所から\qty{100}{MW}の電力を、往復の抵抗が\qty{10}{\ohm}の送電線で送ります（数は仮のものとします）。電圧が\qty{100}{kV}のときと\qty{500}{kV}のときで、電線に流れる電流と、電線で熱になる電力（損失）を求めましょう。損失は、送る電力の何\,\%ですか。
\phbamen
同じ\qty{100}{MW}を、違う電圧で送ります。電圧を高くすると、電線の熱がどう変わるかを知りたい。
\phmushi
本当の送電線は3本1組の三相ですが、行きと帰りの2本（単相）として考えます。電線のコイルやコンデンサーとしての性質（電-13）と、電線での電圧の下がりを無視し、交流を直流と同じように扱います（電流と電圧の山の位置がそろっている、とする）。
\phzu
送る電力 $P=\qty{1.0e8}{W}$、電線の抵抗 $R=\qty{10}{\ohm}$、送る電圧 $V=\qty{1.0e5}{V}$ と\qty{5.0e5}{V}、電流 $I$、損失 $P_{\text{損}}$。
\phriyu
電力の式 $P=VI$（電-4 ⑤）から、同じ電力を送るのに要る電流がわかります。その電流が抵抗 $R$ に流れたときの熱が、ジュール熱 $RI^2$ です。
\phtateru
\[I=\frac{P}{V},\qquad P_{\text{損}}=RI^2\]
\phtoku
\qty{100}{kV}のとき：
\[I=\frac{1.0\times10^8}{1.0\times10^5}=\qty{1000}{A},\qquad P_{\text{損}}=10\times1000^2=\qty{1.0e7}{W}=\qty{10}{MW}\]
\qty{500}{kV}のとき：
\[I=\frac{1.0\times10^8}{5.0\times10^5}=\qty{200}{A},\qquad P_{\text{損}}=10\times200^2=\qty{4.0e5}{W}=\qty{0.40}{MW}\]
\phtashikame
単位：W÷V＝A、Ω×A$^2$＝(V/A)×A$^2$＝V·A＝W。\\
大きさ：\qty{100}{kV}では、送る電力の10\,\%が熱になってしまいます。\qty{500}{kV}なら0.4\,\%です。電圧を5倍にすると、電流は $\tfrac15$、損失は $\tfrac1{25}$ になりました（$10\,\%\div25=0.4\,\%$）。\\
極端な場合：家で使う\qty{100}{V}のまま送ろうとすると、電流は100万 A になり、損失 $10\times(10^6)^2$ W を計算すると、送る電力の10万倍という、ありえない数になります。\qty{100}{V}では、この送電線でとても送れないということです。
\kotae \qty{100}{kV}：\qty{1000}{A}、\qty{10}{MW}（10\,\%）。\qty{500}{kV}：\qty{200}{A}、\qty{0.40}{MW}（0.4\,\%）
\end{phreidai}
\noindent 同じ送電線で、電圧の段ごとに損失をくらべたのが図\ref{ph:電-0:fig-sonshitsu}です。\qty{66}{kV}では2割以上が熱になりますが、\qty{500}{kV}では1\,\%にもなりません。
\begin{bkfigure}
\begin{tikzpicture}[x=1mm, y=1mm]
  % 損失〔%〕を高さ 0.85 mm/% にした棒（目もりの軸は置かず、棒の上に数を書く）
  \draw[bkInk, line width=0.6pt] (0,0) -- (104,0);
  \foreach \x/\h/\v/\t in {14/19.5/66/23, 38/3.59/154/4.2, 62/1.12/275/1.3, 86/0.34/500/0.4} {
    \filldraw[fill=bkRed!45, draw=bkRed!70!black, line width=0.5pt] (\x-7,0) rectangle (\x+7,\h);
    \node[sxa lab, anchor=south] at (\x,\h+0.6) {\t\,\%};
    \node[sxa lab, anchor=north] at (\x,-1.2) {\v\,kV で送る};}
\end{tikzpicture}
\bkfigcaption{\qty{100}{MW}を往復\qty{10}{\ohm}の線で送るときの損失（送る電力の\%）。\bkreidairef{3}と同じ計算。}\label{ph:電-0:fig-sonshitsu}
\end{bkfigure}

\begin{phbaisuru}{送る電圧 $V$}
同じ電力を送るとき、電圧を2倍にすると、要る電流は半分になります（$P=VI$ で $P$ が同じ）。損失 $RI^2$ は電流の2乗なので、$\tfrac12\times\tfrac12=\tfrac14$ になります。\bkreidairef{3}の線を\qty{200}{kV}で送れば、電流は\qty{500}{A}、損失は\qty{10}{MW}の $\tfrac14$ の\qty{2.5}{MW}です。電流を正方形の1辺の長さにすると、損失は正方形の面積のように変わります。
\begin{bkcenter}
\begin{tikzpicture}[x=1mm, y=1mm]
  \filldraw[fill=bkRed!30, draw=bkRed!70!black, line width=0.7pt] (0,0) rectangle (24,24);
  \filldraw[fill=bkRed!30, draw=bkRed!70!black, line width=0.7pt] (40,0) rectangle (52,12);
  \draw[bkRed!70!black, line width=0.4pt, dashed] (40,0) rectangle (64,24);
  \filldraw[fill=bkRed!30, draw=bkRed!70!black, line width=0.7pt] (80,0) rectangle (84.8,4.8);
  \draw[bkRed!70!black, line width=0.4pt, dashed] (80,0) rectangle (104,24);
  \node[sxa lab, align=center, anchor=north] at (12,-1.5) {\qty{100}{kV}\\\qty{1000}{A}・損失 \qty{10}{MW}};
  \node[sxa lab, align=center, anchor=north] at (52,-1.5) {\qty{200}{kV}（2倍）\\\qty{500}{A}・損失 \qty{2.5}{MW}};
  \node[sxa lab, align=center, anchor=north] at (92,-1.5) {\qty{500}{kV}（5倍）\\\qty{200}{A}・損失 \qty{0.4}{MW}};
\end{tikzpicture}
\bkfigcaption{電流（正方形の1辺）と損失（面積）。電圧2倍で面積は $\tfrac14$、5倍で $\tfrac1{25}$。破線は\qty{100}{kV}のときの大きさ。}
\end{bkcenter}
\end{phbaisuru}

\noindent このことを、送る電力 $P$・電圧 $V$・抵抗 $R$ の式にまとめ、力率（電圧と電流の山のずれ）まで入れて確かめるのは、電-11 ⑥です（電-11 の例題5は、この例題と同じ数です）。

\phboxhead{では、電圧はいくらでも高くすればよいのか}
電圧を上げるほど、電線と鉄塔の間、電線どうしの間で、空気がこらえきれずに電気を通してしまう（放電する）心配が大きくなります。そこで、電線を支える部品（がいし）を長くし、鉄塔を高く大きくしなければなりません。家の中で何十万 V を使うこともできません。そのため、遠くへたくさん送る所だけを高い電圧にし、使う所の近くで何段にも分けて下げているのです（図\ref{ph:電-0:fig-tabi}）。電圧を自由に上げ下げできる変圧器が交流でしか働かないことが、電気を交流で送るようになった大きな理由です（電-11 ⑥・電-12 ⑤）。

%% ================================================================
\Needspace{10\baselineskip}
\bksubsection{5}{4つの部と、16のなぜ？}
\leavevmode\label{ph:電-0:sec5}この巻は、序章と4つの部（16の章）でできています。電気の旅の段ごとに、4つの部の物理が働いています。
\begin{itemize}[leftmargin=1.2\zw, itemsep=2pt, topsep=2pt, label={\color{bkTheme}●}]
\item \textbf{第1部　引きあう}：電気の力と電場、電位、コンデンサー、電流と抵抗、回路の計算。ドライヤーの電熱線と、電線の網の計算（電-1〜5）。
\item \textbf{第2部　うず巻く}：電流がつくる磁場、鉄と磁石、電磁誘導。発電機と変圧器が動くしくみの土台（電-6〜9）。
\item \textbf{第3部　変える}：交流、送電の電力と三相交流、モーターと変圧器、長い電線（電-10〜13）。
\item \textbf{第4部　届く}：電気と磁気の4つの式、電磁波（光・電波）、電気を届ける・変える（電-14〜16）。
\end{itemize}
\begin{bkcenter}
\small
\begin{tblr}{colspec={Q[l,m,wd=19mm] X[l,m] Q[l,m,wd=33mm]}, width=\linewidth, rowsep=1.1pt, colsep=3pt,
  hlines={0.4pt,bkLine}, row{1}={bg=bkTheme!12, font=\sffamily\bfseries\footnotesize},
  row{2,8}={bg=bkTheme!6, font=\sffamily\bfseries\footnotesize}}
章 & 身近な「なぜ？」 & 答えのある所 \\
\SetCell[c=3]{l}{\phfaceicon[5mm]{1}\enspace 第1部　引きあう：静電場と電流} & & \\
\phchapref{電-1} & 下じきでこすった髪の毛は、なぜ立ち上がる？ & ② クーロンの法則 \\
\phchapref{電-2} & 電気の力は、なぜはなれるほど弱くなる？ & ④ 対称な電荷 \\
\phchapref{電-3} & 雷が落ちても、金属の車の中の人が無事なのはなぜ？ & ③ 導体の性質 \\
\phchapref{電-4} & ドライヤーは、電気でどうやって熱をつくる？ & ⑤ ジュール熱と電力 \\
\phchapref{電-5} & 鳥は、なぜ電線にとまっても感電しない？ & ① 電流の分かれ方 \\
\SetCell[c=3]{l}{\phfaceicon[5mm]{2}\enspace 第2部　うず巻く：磁場と電磁誘導} & & \\
\phchapref{電-6} & 電線に電気を流すと、なぜ方位磁針がふれる？ & ⑤ まっすぐな電流 \\
\phchapref{電-7} & 磁石にくっついたくぎは、なぜ自分も磁石になる？ & ② 小さな磁石がそろう \\
\phchapref{電-8} & IH調理器は、火がないのになぜなべが熱くなる？ & ⑫ 表皮効果と IH \\
\phchapref{電-9} & 充電器をコンセントからぬいても、ランプがしばらく光っているのはなぜ？ & ⑤ 充電器のコンデンサー \\
\end{tblr}
\end{bkcenter}
\begin{bkcenter}
\small
\begin{tblr}{colspec={Q[l,m,wd=19mm] X[l,m] Q[l,m,wd=33mm]}, width=\linewidth, rowsep=1.1pt, colsep=3pt,
  hlines={0.4pt,bkLine}, row{1,6}={bg=bkTheme!6, font=\sffamily\bfseries\footnotesize}}
\SetCell[c=3]{l}{\phfaceicon[5mm]{3}\enspace 第3部　変える：交流と電気機器} & & \\
\phchapref{電-10} & ラジオは、たくさんの電波からどうやって1つの局を選ぶ？ & ⑦ 共振 \\
\phchapref{電-11} & 送電線の鉄塔では、なぜ電線が3本ずつ組になっている？ & ⑧ 三相の電力 \\
\phchapref{電-12} & 電気自動車のモーターは、どうやって回る？ & ⑧ 同期モーター \\
\phchapref{電-13} & テレビのアンテナ線に、なぜふつうの電線を使わない？ & ⑧ はね返りと整合 \\
\SetCell[c=3]{l}{\phfaceicon[5mm]{4}\enspace 第4部　届く：電磁波} & & \\
\phchapref{電-14} & 電気と磁気は、別々のものなの？ & ⑨ 見る人で入れかわる \\
\phchapref{電-15} & 電子レンジの窓の網は、何のためにある？ & ⑩ 金属と電磁波 \\
\phchapref{電-16} & 停電しないように、電気の量をどうやって合わせている？ & ② 周波数と需給 \\
\end{tblr}
\bktabcaption{この巻の地図。「答えのある所」の丸数字は、その章の「物理の道具」の節の番号。}\label{ph:電-0:tab-chizu}
\end{bkcenter}
表\ref{ph:電-0:tab-chizu}は、それぞれの章で答える身近な「なぜ？」と、その答えがある所です。章の札の横の数は、その章が始まるページです。どの章から読んでもよいのですが、前の章の道具を使うので、はじめは順に読むのがおすすめです。序章の\bkhatenaref の答えは、④と、この序章の最後（⑧）にあります。身近な疑問から、数学と物理のどの章を順に読めばよいかを示した「なぜ？の道すじ」も、巻末にあります。

%% ================================================================
\Needspace{12\baselineskip}
\bksubsection{6}{数学の本とのつながり}
\leavevmode\label{ph:電-0:sec6}電気と磁気の式の多くは、数学の本でもう計算しています。たとえば、コンデンサーの充電は 6-4、交流を回る矢印で表すフェーザは 6-10、電気と磁気の4つの式は 8-14 で扱いました。この巻では、同じ計算を繰り返すかわりに、「（6-10⑨）」のように数学の章と節を指します。そして、物理の本ならではの、次の3つを中心に書きます。
\begin{itemize}[leftmargin=1.2\zw, itemsep=2pt, topsep=2pt, label={\color{bkTheme}●}]
\item その式は、どんな実験と法則から出てくるのか。何を無視して（どんな仮定で）作ったのか。
\item 本物の物（電線・モーター・変圧器・アンテナ）の値を入れると、何が起きるのか。
\item 仮定がくずれると、何が変わるのか。
\end{itemize}
\begin{bkcenter}
\small
\begin{tblr}{colspec={Q[l,m,wd=30mm] X[l,m] Q[l,m,wd=26mm]}, width=\linewidth, rowsep=1.1pt, colsep=3pt,
  hlines={0.4pt,bkLine}, row{1}={bg=bkTheme!12, font=\sffamily\bfseries\footnotesize}}
数学の本の章 & そこで解いたこと & この巻の章 \\
2-15 & 電圧と電流は比例する（オームの法則） & 電-4 \\
8-5・8-6 & 電場の矢印、広がった物の積分 & 電-1 \\
8-8・8-11・8-12 & 面を通りぬける量、ガウスの定理、ラプラシアン & 電-2・電-3 \\
7-5・7-11 & 逆行列で連立方程式を解く、電気回路の行列 & 電-5 \\
4-14・8-10 & 外積と右ねじ、うず（rot） & 電-6・電-12 \\
8-11 & ストークスの定理 & 電-6・電-8 \\
6-4・9-7 & 指数関数で減る量、ラプラス変換 & 電-9 \\
6-9・6-10・6-13 & 複素平面、フェーザ、共振 & 電-10・電-11 \\
7-3・7-4・9-11 & 行列のかけ算と行列式、行列の指数関数 & 電-13 \\
8-14 & 電気と磁気の4つの式、光の速さ & 電-14・電-15 \\
\end{tblr}
\bktabcaption{数学の本で解いた計算と、それを使うこの巻の章。}\label{ph:電-0:tab-math}
\end{bkcenter}
\noindent 数学の本をまだ読んでいなくても、だいじょうぶです。★が5つの内容をふくむ章は、「目でわかる」「手でわかる」「式でわかる」の3段に分けてあり、はじめの2段は、図と表計算を中心に、★3つまでの道具でたどれるように書いてあります。

%% ================================================================
\Needspace{8\baselineskip}
\bksubsection{7}{この本の読み方}
\leavevmode\label{ph:電-0:sec7}
\begin{bkfigure}
\begin{tikzpicture}[x=1mm, y=1mm]
  \sxatalk{0}{sora}{ソラ}{電気って、発電所から電子が家まで走ってくるのかと思っていたよ。}
  \sxatalk{-20}{mina}{ミナ}{わたしは、磁石の話と電波の話が、電気と同じ本にあるのがふしぎ……。}
  \sxatalk{-44}{delta}{デルタ}{電子は1秒に1 mm も進まないよ。届くのは「押せ」という知らせで、エネルギーは電線の周りの電場と磁場が運ぶんだ（電-14）。電気と磁気は1つのものだから、1冊の本になるんだね。★が5つの所は3段に分けてあるから、「ここまでで十分」の旗まで読めば、その章のゴールだよ。}
\end{tikzpicture}
\end{bkfigure}
\begin{itemize}[leftmargin=1.2\zw, itemsep=2pt, topsep=2pt, label={\color{bkTheme}●}]
\item \textbf{★の数}：内容の高さの目安です。★1つは「はじめの一歩」、★2つは「ふもとを歩く」、★3つは「山を登る」、★4つは「尾根を歩く」、★5つは「頂上をめざす」。物理の本の★は、物理の内容の高さで付けてあり、数学の本の★とは別の物差しです。
\item \textbf{章の札}：「8-14」のような札は数学の本の章、「電-4」のような札はこの本の章、「力-4」「熱-0」「波-7」は第1巻『力学』・第2巻『熱学』・第3巻『波動』の章です。札の色は、その本の色です。本文で「（6-10⑨）」と書いたら、数学の本の6-10の⑨の節を見てください。
\item \textbf{物理の本の囲み}：「この式の読み方」「身近な例では」「〇〇を2倍にすると？」「まだわかっていないこと」「この先の学び」「一般の式から、よく使う式へ」などがあります。見本は、前付の「この本の使い方」にあります。
\item \textbf{本物の値}：銅の抵抗率や、家の電圧の決まりなど、この巻で使う値は、巻末の「物性値・定数・単位」にまとめてあります。
\end{itemize}
\begin{taisetsu}
\def\tsl#1{\par\noindent\hangindent=6.2\zw\hangafter=1\makebox[6.2\zw][l]{\sffamily\bfseries\small\color{bkDeep}#1}}
\tsl{4つの量}電圧 $V$〔V〕、電流 $I$〔A〕、抵抗 $R$〔\si{\ohm}〕、電力 $P$〔W〕。$V=RI$、$P=VI$。
\tsl{電位と電圧}場の章では、場所ごとの電気の「高さ」（電位）を $\phi$、2点の電位の差（電圧）を $V$ と書く（電-2）。
\tsl{回路の記号}回路の章（電-9〜13）では、電源の電圧を $e$・$E$、時間で変わる電圧・電流を小文字の $v$・$i$ と書く。
\tsl{交流}交流の電圧・電流は、ことわりがなければ実効値。「\qty{100}{V}」の山は約\qty{141}{V}。
\tsl{答えのけた}答えは、ふつう2けた。数学の本と同じ数を使う例題などは3けたで書く。
\end{taisetsu}
\noindent 記号の読み方と意味は、巻末の「記号カード」にまとめてあります。

%% ================================================================
\Needspace{6\baselineskip}
\bksubsection{8}{はてな？の答え：電流を小さくして送る}
\leavevmode\label{ph:電-0:sec8}\bkhatenaref に答えましょう。
\begin{enumerate}[label={\bkcircled[bkBlue]{\arabic*}}, leftmargin=2.4\zw, itemsep=2pt, topsep=2pt]
\item \textbf{電線は熱を出す}　送電線にも抵抗があり、電流 $I$ が流れると、$RI^2$ の電力が熱になって失われます。失われる電力は、\textbf{電流の2乗}に比例します（④、電-4 ⑤）。
\item \textbf{電圧を上げると、電流が減る}　送る電力は $P=VI$ です。同じ電力を送るなら、電圧を10倍にすれば、電流は $\tfrac1{10}$ ですみます（②）。
\item \textbf{だから、損失は2乗で減る}　電流が $\tfrac1{10}$ なら、損失は $\tfrac1{100}$ です。\qty{100}{MW}を送るとき、\qty{100}{kV}では10\,\%が熱になりますが、\qty{500}{kV}なら0.4\,\%ですみます（\bkreidairef{3}）。式で確かめるのは電-11 ⑥です。
\item \textbf{上げ下げは変圧器で}　交流なら、変圧器で電圧をかんたんに上げ下げできます（電-12 ⑤）。遠くへ送る所だけ電圧を高くし、家の近くで\qty{100}{V}・\qty{200}{V}に下げます（①）。
\end{enumerate}
つまり、送電線がとても高い電圧で電気を送るのは、\textbf{電流を小さくして、電線で熱になって失われる電力を減らすため}です。
この電気の旅には、4つの部の物理が、すべて働いています（図\ref{ph:電-0:fig-kotae}）。
\begin{bkfigure}
\begin{tikzpicture}[x=1mm, y=1mm]
  \tikzset{tag/.style={draw=#1, line width=0.9pt, fill=white, rounded corners=1.6mm, text width=58mm, align=left, inner sep=1.8mm, font=\sffamily\footnotesize}}
  \node[tag=bkBlue, anchor=north west] (n1) at (0,0) {\phfaceicon[5mm]{1}\enspace\textbf{引きあう}\\電線の中で電場が電子を押し、電流が流れる。電線の抵抗で、電力の一部が熱になる。\\\phchapref{電-4}\enspace\phchapref{電-5}};
  \node[tag=bkOrange!80!black, anchor=north west] (n2) at (66,0) {\phfaceicon[5mm]{2}\enspace\textbf{うず巻く}\\発電機は、磁場の変化で電圧を生む。変圧器も同じ電磁誘導で働く。\\\phchapref{電-6}\enspace\phchapref{電-8}};
  \node[tag=bkGreen!70!black, anchor=north west] (n3) at (0,-35) {\phfaceicon[5mm]{3}\enspace\textbf{変える}\\交流は、変圧器で電圧を上げ下げできる。送電線は三相で送る。\\\phchapref{電-11}\enspace\phchapref{電-12}};
  \node[tag=bkPurple, anchor=north west] (n4) at (66,-35) {\phfaceicon[5mm]{4}\enspace\textbf{届く}\\電気のエネルギーは、電線の周りの電場と磁場が、光の速さに近い速さで運ぶ。\\\phchapref{電-13}\enspace\phchapref{電-14}};
\end{tikzpicture}
\bkfigcaption{発電所から家までの電気の旅に働く、4つの部の物理。}\label{ph:電-0:fig-kotae}
\end{bkfigure}
\begin{column}{この本の「やってみよう」の約束}
家のコンセントの\qty{100}{V}や、それより高い電圧は、体に流れると命にかかわります（電-4 ⑫）。そこで、この本の「やってみよう」では、乾電池・磁石・アルミはく・方位磁針・スマホのアプリまでしか使いません。電子レンジ・IH調理器・送電線は、「見るだけ・読むだけ」にします。乾電池でも、＋と−を導線だけで直接つなぐと、導線と電池が熱くなるので、必ず豆電球などを入れてください。電柱や鉄塔に登ったり、切れて垂れた電線に近づいたりしてはいけません。電気のこわさを正しく知ることも、電磁気学の大切な一歩です。
\end{column}

\begin{phketa}[電気の旅の数]
\phketaline{乾電池・家のコンセント}{\qty{1.5}{V}・\qty{100}{V}（\qtyrange{95}{107}{V}に保つ）と\qty{200}{V}}
\phketaline{配電線・送電線}{\qty{6600}{V}・\qty{66}{kV}〜\qty{500}{kV}}
\phketaline{ドライヤーの電流と電力}{\qty{12}{A}・\qty{1200}{W}（\bkreidairef{1}）}
\phketaline{\qty{100}{MW}を\qty{500}{kV}で送る電流}{\qty{200}{A}（\bkreidairef{3}）}
\phketaline{電線の中の電子の速さ}{1秒に約\qty{0.44}{mm}（電-4 例題2）}
\phketaline{電気の知らせが\qty{400}{km}進む時間}{約\qty{1.3}{ms}（\bkreidairef{2}）}
\phboxhead{答えが桁で合っているか}
家の電流は数 A〜十数 A、送電線の電流は数百〜千 A くらいです。送電の損失が送る電力をこえたり、電子が光の速さで走ったりしたら、kV と V、MW と W を取り違えているか、電子の速さと知らせの速さを取り違えています。
\end{phketa}

\begin{phmada}[電気を、たくさん、長くためておくには？]
発電所でつくる電気の量と、みんなが使う電気の量は、いつも同じにしておかなければなりません。電気は、つくったそばから使われていくからです。どうやって合わせているかは、電-16 で学びます。もし、電気を安く、たくさん、長くためておくことができれば、太陽の光や風でつくった電気を、夜や風のない日にも使えます。いまは、電気が余ったときに水を高い所へくみ上げておき、足りないときに流して発電する\textbf{揚水発電}や、大きな蓄電池が使われています。\bkidx{ようすいはつでん}{揚水発電}でも、何日も何か月もの間、たくさんの電気を安くためておく方法は、まだ研究の途中です。たとえばアメリカのエネルギー省は、2021年に、10時間をこえて電気をためる大きな蓄電の費用を10年で9割下げることを目標にかかげ、さまざまな方法の研究を進めています。どの方法がいちばんよいかは、まだわかっていません（電-16）。
\end{phmada}


\begin{kotoba}\small
\bkword{回路}{電気が流れる、ひとまわりの輪の道すじ。}
\bkword{交流・直流}{向きと大きさが周期的に変わる電気が交流、いつも同じ向きの電気が直流。}
\bkword{変圧器}{2つのコイルの巻数の比で、交流の電圧を上げたり下げたりする機械。}
\bkword{送電線・配電線}{発電所から町へ高い電圧で送る線と、町の中で配る\qty{6600}{V}などの線。}
\bkword{送電の損失}{電線の抵抗で熱になって失われる電力。電流の2乗に比例する。}
\end{kotoba}

\begin{furikaeri}
\begin{checklist}
\item 発電所から家のドライヤーまで、電気が通る道すじと、とちゅうで働く機械（発電機・変圧器・送電線）を説明できる。\bkfurref{ph:電-0:sec1}
\item 何を無視したかを書いてから、電流・電力・電線で失われる電力を、7つの段で求められる。\bkfurref{ph:電-0:ex3}
\item 身近な「なぜ？」の答えが、この巻のどの章にあるかを探せる。\bkfurref{ph:電-0:sec5}
\end{checklist}
\end{furikaeri}
\bkfukidashi{sora}{高い電圧で送るのは、電流を小さくするためだったんだね。まずは電-1 で、電気の力そのものから始めよう。}
